\section*{Bab \ref{ch:b3:solid-state} --- Kimia Zat Padat: Pita, Cacat, dan Semikonduktor}

\begin{solution}{exo:b3:solid-state:1}
$(E - \alpha)/|\beta| = -2\cos(k\pi/7)$: $-1.802$, $-1.247$, $-0.445$, $+0.445$, $+1.247$, $+1.802$. Tingkat-tingkat itu mencakup
$3.604|\beta|$, sudah 90\,\% dari lebar pita $4|\beta|$ rantai tak hingga.
\end{solution}

\begin{solution}{exo:b3:solid-state:2}
$2N$ ($N$ orbital, masing-masing dua elektron). Natrium, dengan satu elektron 3s per atom, mengisi
setengah pita s-nya: sebuah logam. Magnesium memiliki dua elektron dan akan mengisi pita s tepat penuh,
tetapi pita 3s dan 3p bertumpang tindih, sehingga tingkat-tingkat terisi tertinggi termasuk dalam pita
gabungan yang terisi sebagian: magnesium juga logam.
\end{solution}

\begin{solution}{exo:b3:solid-state:3}
$\lambda \approx \qty{1239.8}{eV.nm}/E_g$: galium fosfida \qty{549}{nm} (hijau); galium nitrida \qty{365}{nm}
(ultraviolet dekat). Dioda biru memakai galium nitrida yang dipadu dengan indium, yang
celahnya lebih kecil.
\end{solution}

\begin{solution}{exo:b3:solid-state:4}
$\ce{CaCl2} \xrightarrow{\ce{NaCl}} \mathrm{Ca_{Na}^{\bullet} + V_{Na}' + 2\,Cl_{Cl}^{\times}}$: dua situs kation dan dua situs anion, seperti dalam NaCl; satu Ca dan
dua Cl; muatan $+1 - 1 = 0$. $\ce{Y2O3} \xrightarrow{\ce{ZrO2}} \mathrm{2\,Y_{Zr}' + 3\,O_O^{\times} + V_O^{\bullet\bullet}}$: dua situs kation dan empat situs oksigen; dua Y
dan tiga O; muatan $-2 + 2 = 0$.
\end{solution}

\begin{solution}{exo:b3:solid-state:5}
$T^{3/2} = 5196$ pada \qty{300}{K}, $kT = \qty{0.02585}{eV}$. Silikon: $\sqrt{N_cN_v} = \qty{2.42e19}{cm^{-3}}$, $\eu^{-1.125/0.05170} = \num{3.55e-10}$,
$n_i = \qty{8.6e9}{cm^{-3}}$. Germanium: $\sqrt{N_cN_v} = \qty{7.16e18}{cm^{-3}}$, $\eu^{-0.661/0.05170} = \num{2.80e-6}$, $n_i = \qty{2.0e13}{cm^{-3}}$. Rasionya
$\num{4.3e-4}$: celah germanium yang lebih kecil memberinya lebih dari dua ribu kali lebih banyak
\omterm{def:b3:solid-state:carriers}{pembawa muatan}.
\end{solution}

\begin{solution}{exo:b3:solid-state:6}
$n = N_D = \qty{1.0e16}{cm^{-3}}$; $p = n_i^2/n = \num{1.0e20}/\num{1.0e16} = \qty{1.0e4}{cm^{-3}}$.
\end{solution}

\begin{solution}{exo:b3:solid-state:7}
$E_F - E_i = kT\ln(n/n_i) = \qty{0.02585}{eV} \times \ln(10^6) = \qty{0.36}{eV}$.
\end{solution}

\begin{solution}{exo:b3:solid-state:8}
$kT = \qty{0.06894}{eV}$ pada \qty{800}{K}; $n/N = \eu^{-2.3/0.1379} = \num{5.7e-8}$; dalam satu mol $\num{6.02e23} \times \num{5.7e-8} = \num{3.4e16}$
pasangan.
\end{solution}

\begin{solution}{exo:b3:solid-state:9}
Massa satu sel adalah $\rho a^3 = \qty{5.70}{g.cm^{-3}} \times (\qty{4.300e-8}{cm})^3 = \qty{4.532e-22}{g}$, yaitu
$\qty{4.532e-22}{g} \times N_A = \qty{272.9}{g}$ per mol sel, atau \qty{68.23}{g} per oksigen. Maka
$(1 - x) \times 55.8 + 16.0 = 68.23$ memberikan $1 - x = 0.936$: $x = 0.064$, $\mathrm{Fe_{0.936}O}$.
\end{solution}

\begin{solution}{exo:b3:solid-state:10}
$\ln(\sigma T)$: $-3.51$, $-1.30$, $0.36$ pada $1/T = 3.333$, $2.857$, $\qty{2.500e-3}{K^{-1}}$; kemiringan antara kedua titik ujung
$(0.36 + 3.51)/(-\qty{8.33e-4}{K^{-1}}) = -\qty{4.65e3}{K}$ (titik tengahnya terletak pada garis itu), sehingga
$E_a = \qty{4.65e3}{K} \times \qty{8.617e-5}{eV.K^{-1}} = \qty{0.40}{eV}$.
\end{solution}

\begin{solution}{exo:b3:solid-state:11}
$\tfrac12\ce{O2} \rightleftharpoons \mathrm{O_O^{\times} + V_M' + h^{\bullet}}$, $K = [\mathrm{V_M'}][\mathrm h^\bullet]/p(\ce{O2})^{1/2}$. Neraca muatan $[\mathrm h^\bullet] = [\mathrm{V_M'}]$, sehingga
$[\mathrm h^\bullet]^2 = K\,p(\ce{O2})^{1/2}$ dan $[\mathrm h^\bullet] \propto p(\ce{O2})^{1/4}$; konduktivitasnya mengikuti \omterm{def:b3:solid-state:carriers}{lubang}.
\end{solution}

\begin{solution}{exo:b3:solid-state:12}
$n/N = \sqrt{N_i/N}\,\eu^{-\Delta H_F/2kT} = \sqrt 2\,\eu^{-1.4/(2 \times 0.05170)} = 1.414 \times \num{1.32e-6} = \num{1.9e-6}$.
\end{solution}

\begin{solution}{pb:b3:solid-state:1}
\textbf{1.} $\ce{Y2O3} \xrightarrow{\ce{ZrO2}} \mathrm{2\,Y_{Zr}' + 3\,O_O^{\times} + V_O^{\bullet\bullet}}$.

\textbf{2.} Situs: dua situs kation dan empat situs oksigen (tiga terisi, satu kosong),
yaitu rasio 1\,:\,2 pada $\ce{ZrO2}$. Massa: 2 Y dan 3 O di setiap sisi. Muatan: $2 \times (-1) + 2 = 0$.

\textbf{3.} Satu kekosongan untuk dua ion itrium: setengah per itrium.

\textbf{4.} Per $(\ce{ZrO2})_{0.92}(\ce{Y2O3})_{0.08}$ terdapat 0.92 Zr dan 0.16 Y, yaitu 1.08 kation: $x = 0.16/1.08 = 0.148$,
$\mathrm{Zr_{0.852}Y_{0.148}O_{1.926}}$.

\textbf{5.} $(x/2)/2 = 0.0370$: 3.7\,\% situs oksigen kosong.

\textbf{6.} $8 \times 0.0370 = 0.296$ kekosongan per sel.

\textbf{7.} $a^3 = (\qty{5.14e-8}{cm})^3 = \qty{1.358e-22}{cm^3}$: $0.296/\qty{1.358e-22}{cm^3} = \qty{2.2e21}{cm^{-3}}$.

\textbf{8.} Setiap ion oksida harus melompati penghalang \qty{1.0}{eV} ke dalam kekosongan
yang bertetangga; fraksi percobaan yang berhasil, $\eu^{-E_a/kT}$, tumbuh sangat cepat bersama $T$.

\textbf{9.} Pada \qty{1073.15}{K}, $kT = \qty{0.09248}{eV}$: $A = \sigma T\eu^{E_a/kT} = 0.030 \times 1073.15 \times \eu^{10.81} = \qty{1.6e6}{S.K.cm^{-1}}$.

\textbf{10.} Pada \qty{973.15}{K}: $\sigma = (A/T)\eu^{-11.92} = \qty{0.011}{S.cm^{-1}}$.

\textbf{11.} Pada \qty{573.15}{K}: $\sigma = (A/T)\eu^{-20.25} = \qty{4.5e-6}{S.cm^{-1}}$.

\textbf{12.} $R = L/(\sigma S) = \qty{0.10}{cm}/(\sigma \times \qty{0.20}{cm^2})$: \qty{46}{\ohm} pada \qty{700}{\celsius}, $\qty{1.1e5}{\ohm}$ pada \qty{300}{\celsius}.

\textbf{13.} $R < \qty{1}{k\ohm}$ memerlukan $\sigma > \qty{5.0e-4}{S.cm^{-1}}$; menyelesaikan $(A/T)\eu^{-E_a/kT} = \num{5.0e-4}$ (dengan coba-coba)
memberikan $T \approx \qty{760}{K}$, sekitar \qty{490}{\celsius}.

\textbf{14.} Gas buang masih dingin setelah mesin dihidupkan dan pada beban rendah; pemanas membawa
elektrolit ke suhu kerjanya dalam beberapa detik, sehingga sistem kendali
bekerja sejak awal, saat sebagian besar polutan diemisikan.

\textbf{15.} $\ce{O2 + 4e- -> 2O^2-}$; dalam \omterm{def:b3:solid-state:kroger-vink}{notasi Kröger--Vink} $\ce{O2} + \mathrm{2\,V_O^{\bullet\bullet} + 4\,e'} \longrightarrow \mathrm{2\,O_O^{\times}}$.

\textbf{16.} Oksigen direduksi pada sisi udara dan ion-ion oksida dioksidasi kembali
menjadi oksigen pada sisi gas buang: sel itu memindahkan $\ce{O2}$ dari udara ke gas buang, dengan
$\Delta_rG = RT\ln(p_{\mathrm{buang}}/p_{\mathrm{udara}})$ per mol $\ce{O2}$ dan empat elektron: $E = -\Delta_rG/4F = (RT/4F)\ln(p_{\mathrm{udara}}/p_{\mathrm{buang}})$.

\textbf{17.} $RT/4F = 8.3145 \times 973.15/(4 \times 96485) = \qty{0.02096}{V}$.

\textbf{18.} Nol: kedua tekanan itu sama.

\textbf{19.} $(RT/4F)\ln 10 = \qty{48.3}{mV}$ per dekade.

\textbf{20.} $\qty{0.02096}{V} \times \ln(0.2095/0.010) = \qty{0.02096}{V} \times 3.04 = \qty{0.064}{V}$.

\textbf{21.} $\qty{0.02096}{V} \times \ln(0.2095/\num{1.0e-20}) = \qty{0.02096}{V} \times 44.49 = \qty{0.93}{V}$.

\textbf{22.} Gas buang kaya mengandung karbon monoksida, hidrogen, dan hidrokarbon
yang tidak terbakar; pada elektrode platina, sedikit oksigen yang tersisa bereaksi dengan zat-zat itu,
dan $p(\ce{O2})$ ditentukan oleh kesetimbangan $\ce{2CO + O2 <=> 2CO2}$ dan $\ce{2H2 + O2 <=> 2H2O}$, yang menyisakan nilai yang
sangat kecil pada \qty{700}{\celsius}.

\textbf{23.} $p = 0.2095\,\eu^{-0.45/0.02096} = \qty{1.0e-10}{bar}$.

\textbf{24.} Di sekitar rasio stoikiometrik, gas buang berubah dari sedikit kelebihan
oksigen menjadi sedikit kelebihan CO dan hidrogen: $p(\ce{O2})$ turun sekitar delapan belas
orde besaran untuk perubahan bahan bakar yang kecil, dan tegangannya melonjak dari di bawah
\qty{0.1}{V} menjadi sekitar \qty{0.9}{V}. Sistem kendali mesin menambah bahan bakar ketika tegangannya rendah dan
menguranginya ketika tinggi, sehingga campuran tetap pada titik stoikiometrik, tempat
katalis tiga arah mengonversi CO, hidrokarbon, dan nitrogen oksida sekaligus.

\textbf{25.} Dalam gas buang kaya pada \qty{700}{\celsius}, sensor itu memberikan sekitar \qty{0.93}{V}.
\end{solution}
